\documentclass[12pt]{article}

  \usepackage[a4paper,left=20mm, right=25mm,top=25mm, bottom=25mm]{geometry}
  \usepackage[utf8]{inputenc}
  \usepackage[ngerman]{babel}
  \usepackage[T1]{fontenc}

  %Grafiken
  \usepackage{graphicx}
 
  %Tabellen
  \usepackage{array}
  \usepackage{tabularx}
  \usepackage{booktabs}
  \newcolumntype{K}[1]{S[table-format={#1}, table-alignment=center, table-number-alignment=center, input-decimal-markers={,}, output-decimal-marker={,}]}
  \newcolumntype{L}[1]{>{\raggedright\arraybackslash}p{#1}} % linksbündig mit Breitenangabe  
  \newcolumntype{C}[1]{>{\centering\arraybackslash}m{#1}} % zentriert mit Breitenangabe  
  \newcolumntype{R}[1]{>{\raggedleft\arraybackslash}p{#1}} % rechtsbündig mit Breitenangabe  
  
  %Matheformeln
  \usepackage{amsmath,amsfonts,amssymb,amsthm} %mathe symbole
    \newtheorem{thm}{Satz} %[section]
    \newtheorem{defn}[thm]{Definition}
    \newtheorem{lem}[thm]{Lemma}
    \newtheorem{prop}[thm]{Proposition}
    \newtheorem{cor}[thm]{Korollar}

\usepackage{hyperref}

\usepackage{mathptmx}
\usepackage[scaled=.90]{helvet}
%\usepackage{courier}
    
 
%***********************************************************************
\begin{document}

\section{Das Problem}
Es sei $n\in\mathbb{N}$ mit $n\geq 2$. Wir definieren die Indexmenge $N=\{1, 2,\dots, n\}$ und zwei Vektoren $p = (p_1, p_2, \dots, p_n), q = (q_1, q_2, \dots, q_n)$ mit $n\geq 2$ und $p_i, q_i \in \{0, 1\}$ für $i\in N$. 
Für $i,j \in N$ mit $i<j$ sei 
\begin{align*}
\mathrm{one}_i^j(p) &=   \#\{ k \mid i< k < j, p_k=1 \}
\\
\overline{\mathrm{one}}_i^j(p) &=  \#\{ k \mid 1\leq k<i, p_k=1 \}
+
\#\{ k \mid j< k\leq n, p_k=1 \}
\\
\mathrm{zero}_i^j(p) 
&=   \#\{ k\mid i< k < j, p_k=0 \}
\\
\overline{\mathrm{zero}}_i^j(p) 
&=   \#\{ k\mid 1\leq k < i, p_k=0 \}
+  \#\{ k\mid j< k \leq n, p_k=0 \}
\end{align*}

\begin{defn}
Es sei das Tupel $(n, p, q)$ wie oben beschrieben.
Gibt es $i, j \in N$ mit $i<j$, $p_i = p_j$ und
\begin{align}
\mathrm{one}_i^{j}(p)
&> 
\overline{\mathrm{one}}_i^{j}(p),
\label{eq:Voraussetzung für Tausch -> ones_greater}
\\
\mathrm{zero}_i^{j}(p)
&> 
\overline{\mathrm{zero}}_i^{j}(p),
\label{eq:Voraussetzung für Tausch -> zeros_greater}
\end{align}
dann kann für $(i, j)$ folgende Transformation durchgeführt werden, die wir Tausch oder Wechsel nennen
\begin{align*}
\tau_{ij}(p) 
&= (p_1, p_2, \dots, p_{i-1}, 1-p_i, p_{i+1},\dots, p_{j-1}, 1-p_j, p_{j+1},\dots, p_n).
\end{align*}
Kann ein $q$ durch eine endliche Folge von $m$ Wechseln aus $p$ erreicht werden, also
\begin{align}
(\tau_{i_mj_m} \circ \cdots \circ \tau_{i_2j_2}\circ \tau_{i_1j_1}) (p) 
&= 
q,
\label{eq:p=q durch wechsel}
\end{align}
schreiben wir auch $p\sim q$.
\end{defn}

\begin{defn}
Es sei $(n, p, q)$ gegeben mit $\#\{ i \in N\mid p_i=1 \}$ und $\#\{ i \in N\mid p_i=0 \}$ ungerade.
Gilt $p\sim q$, so nennen wir $(n, p, q)$ lösbar.
\label{defn:problembeschreibung}
\end{defn}
\noindent Und die Frage ist natürlich: Es sei das Tupel $(n, p, q)$ gegeben. Ist $(n, p, q)$ lösbar?

\end{document}